\documentclass{beamer}

\usetheme{Boadilla}
\definecolor{color1}{RGB}{30,50,255}
\definecolor{color2}{RGB}{210,80,40}
%\setbeamercolor{background canvas}{bg=black}
\setbeamercolor{normal text}{fg=color1}
\setbeamercolor{structure}{fg=color1}




\title{Computer TA Class}
\author{University of Tehran}


\date{Autumn 2026}

\begin{document}
	\frame{\titlepage}
		\begin{frame}{Contents}
		
		\begin{itemize}	
			
			\item Scientific calculator $\rightarrow$ first session
			\item Introduction to computer/OS/Bin/Hex $\rightarrow$ first session
			\item Calculations of Octave/Matlab $\rightarrow$ 3-4 sessions
			\item Programming via Octave/Matlab $\rightarrow$ rest of sessions
			
		\end{itemize}
	\end{frame}
	
	\begin{frame}{Software}
		\[
			what\; to\; install
			\begin{cases}
				Windows & Octave/Matlab \\
				Linux/MacOS/BSD & Octave
			\end{cases}
		\]
	\end{frame}
	
	\begin{frame}{Octave}
		Octave is an open-source and free alternative for Matlab made by GNU. No to piracy! 
		\vspace*{1cm}
		\[
		Octave
		\begin{cases}
			Calculation  \\
			programming  
		\end{cases}
		\]
		We are gonna explore the \textbf{Documentation} of the program.
	\end{frame}
	
	\begin{frame}{Calculation}
		\begin{itemize}
			\item Basic calculations and defining matrices and vectors.
			\item More complex calculations and handling experimental data
			\item Statistical calculations
			\item 2D and 3D plotting and some more
			\item equation solving and symbolic package
		\end{itemize}
	\end{frame}
	
	\begin{frame}
		\begin{itemize}
			\item Hardware
			\[
				Essential\; Components \begin{cases}
					??\\ \vdots \\ ??
				\end{cases}
			\]
			\item Software  \[OS \begin{cases}
				DOS/Windows \rightarrow 11,10,XP,98,95,3.1, FreeDOS,\dots \\
				Unixlike \begin{cases}
					BSD \rightarrow OpenBSD, FreeBSD, iOS, macOS, \dots \\
					Linux \rightarrow Debian, Fedora, redhat,  android, \dots
				\end{cases}
			\end{cases} \]
			\item Base-n/logic gates
		\end{itemize}
	\end{frame}
	
	\begin{frame}{Calculators}
		\[
		Calculators
		\begin{cases}
			Basic  \\
			Scientific
				\begin{cases}
					programmble\rightarrow fx-5800, fx-4500, fx-6300, \dots \\
					non-programmble \rightarrow fx-991, fx-82, \dots \\
					Graphical \rightarrow fx-7000,CG-100, \dots
				\end{cases}  
		\end{cases}
		\]
	\end{frame}
	
	\begin{frame}
		\Large
		"help" command and the Documentation are your best tools.
	\end{frame}
	
	\begin{frame}{Octave-start}
		using octave as a calculator...
		\begin{itemize}
			\item data types (int, float, long, string, char, \dots)
			\item simple operations
			\item inf,pi, intmax, realmax,\dots
			\item defining variables(be aware of key words...) "iskeyword('foo')" and "isvarname('foo')"
			\item use "," to combine several statements
			\item using clear foo , clear all, clc , \dots
			\item using whos, who
			\item i, j and complex numbers(real, imag, abs,conj)
			\item sin,sind, sinh, asin,\dots
			\item sqrt, nthroot, factorial, log, log10, gcd, lcm, rem, mod, \dots
			\item round, fix, floor, ceil
			
		\end{itemize}
	\end{frame}
	
	\begin{frame}{vectors \& matrices}
		\begin{itemize}
			\item defining: simple version and n:step:m
			\item ones, zeros, eye,\dots
			\item linespace, logspace
			\item rand, randn, randi, \dots
			\item extraction and deletion from a vector or a matrix
			\item inv, det, transpose, max, mean, sort, diag, sum, prod
			\item addition, subtraction, division and multiplication of elements.
			
			\[
			division
				\begin{cases}
					XA=B \rightarrow X=B/A \\
					AX=B \rightarrow X=A\backslash B
					
				\end{cases}
			\]
		\end{itemize}
		\end{frame}
		
		\begin{frame}{Some other commands}
			\begin{itemize}
				\item dot, cross
				\item xor,and, all, any, find
				\item dec2bin, dec2hex,\dots
				
			\end{itemize}
		\end{frame}
		
		\begin{frame}{Functions}
			\large function [output arguments] = function\_name(input arguments)
			example:\\
			
		\end{frame}
		
		\begin{frame}{Anonymous functions}
			\large f = @(arguments) 'expression'\\
			\vspace{1cm}
			example:\\
			$$y = @(x)\; x.^2+2*x+3;$$
		\end{frame}
		
		\begin{frame}{Polynomial}
				\begin{enumerate}
					\item Define an array
					\item Use '$polyval(P,)$' command to get the value
					\item Use '$roots(P)$' to get the all real and complex roots
					\item '$conv $' and '$ deconv$' commands
					\item $polyfit(y,x,n)$ command
				\end{enumerate}
		\end{frame}
\end{document}